% Page 012 — page imprimée 8
% Fin de l'interrogatoire de Jacques Crégut
% Reconstruction éditable en LaTeX.


Bréau qui se dit cousin du Sr. Delenne, Delenne les aurait quitté et s’en serait allé avec son cousin vers Bréau.

Le répondant et le Sr. Martin étant allés à Aulas, étant arrivés environ 9h00 du soir.

Le répondant connaissait la servante d’Alibert qui est de son village, il serait passé par la porte de derrière de la maison du Sr. Alibert et serait entré dans la chambre de derrière où ils auraient trouvé des personnes d’Aulas que le répondant ne connaît pas qui étaient en débauche... La servante nous ayant donné à manger et à boire sans qu’il y ait Alibert ni sa femme.

Le répondant portait son justaucorps de minime, de même couleur que la veste, son chapeau bordé ayant à la base (?) un ruban vert et un ruban bleu à la chemise.

Et dans le temps qu’il soupait chez le Sr. Alibert ils ont été surpris par la troupe de la garnison de la présente ville, lesquelles ont tué sur le champ le Sr. Martin, les autres qu’il trouva en entrant dans la maison, Alibert et sa femme. Le répondant ayant demandé la vie en disant qu’il découvrirait ses complices on ne l’a pas tué mais on l’a arrêté et mené aux présentes prisons, lui ayant ôté son justaucorps et son chapeau et lui en ayant baillé un autre et lui ayant pris les balles et la poudre qu’il portait ensemble, 3 livres et 4 sols en argent qu’il avait sur lui et c’est le sujet de sa détention.

\vspace{0.35em}

\textit{Interrogé sur ce qu’il était à la St. Michel dernier ?}

Il a été à son village depuis le mois de mai 1702, et 5 à 10 jours il fut à Camprieu pour dépiquer le blé de sa sœur.

\vspace{0.2em}

\textit{Interrogé s’il portait un surnom ?}

Il répond qu’on le nommait Fouzet par sobriquet.

\vspace{0.2em}

\textit{Interrogé s’il avait été au Château du Rey ?}

Il répond que non.

\vspace{0.2em}

\textit{Interrogé s’il avait été à la Lusette avec Teulon ?}

Il répond que non.

\vspace{0.2em}

\textit{Interrogé s’il connaissait Castanet avant Saumane ?}

Il répond que non.

\vspace{0.65em}

\begin{tcolorbox}[
  colback=white,
  colframe=black,
  boxrule=0.7pt,
  arc=0pt,
  left=0.25cm,
  right=0.25cm,
  top=0.12cm,
  bottom=0.12cm,
  boxsep=0pt
]
Après cet interrogatoire on perd la trace de Jacques Crégut

\vspace{0.35em}

\textit{Il est probable qu’un engagement dans l’armée lui permit d’échapper aux poursuites et le récompensa des renseignements donnés : cela dut lui permettre également d’échapper à la vengeance prévisible de ses anciens compagnons}
\end{tcolorbox}

